\cotefile{137}
\scanpage{45}

% 1–3, upper edge
produits $\prod_{i}$ aussi -- ce qui est \deleted{\unsure{vrai|même}} automatique si les $K^i$ sont nuls sauf un nb fini…]

% 4–6, upper struck paragraph
\deleted{Comme le foncteur $\operatorname{Hom}^{\bullet}_{k}(K^{\bullet},C^{\bullet})$ $(B^\circ\longrightarrow(Ab))$ dépend de façon additive des $K^i$, sa représentabilité par $K^{\bullet}$}

% 7
Définissons un hom dans $\widehat B$

% 8, (1)
\[
 F(x)\longrightarrow\operatorname{Hom}^{\bullet}_{k}(\unsure{K^{\bullet}(x)|\alpha^{\bullet}(x)},C^{\bullet}_{F})=\operatorname{Hom}^{\bullet}_{k}(K^{\bullet}(x),F(C^{\bullet})) \tag*{(1)}
\]

% 9
i.e. un hom dans $B^{\bullet}$

% 10, (1′), with brace and gloss
\[
 \underbrace{\deleted{\unsure{x}}F(x)\otimes K^{\bullet}(x)}_{\text{\unsure{représ.|repère}\ \textquotedblright}}\longrightarrow C^{\bullet}=F(C^{\bullet}) \tag*{(1$'$)}
\]

% 10, left margin, leader to underbrace
\margin{\unsure{répr.|rep.} car les $K^i(x)$ sont libres de r.f.}

% 11, (1″), with lower annotations
\[
 \underset{\substack{\rotatebox{-90}{$\in$}\\ B}}{F(x)}\otimes
 \underset{\substack{\rotatebox{-90}{$\in$}\\ Ab_k^{\bullet}}}{\operatorname{Hom}(x,C^{\bullet})}
 \underset{\mathrm{can}}{\longrightarrow}F(C^{\bullet}) \tag*{(1$''$)}
\]

% 12–13
qui provient de l'hom fonctoriel en $\deleted{\unsure{K|x}}C\in\operatorname{Ob}K$

% 14
\[
 F(x)\otimes\operatorname{Hom}(x,C)\longrightarrow F(C)
\]

% 15–20
Comme (1$''$), donc (1) est fonctoriel en $K$, et les deux membres sont des foncteurs additifs en $K$, pour prouver que \deleted{\unsure{c'est}} (1) est un iso, $\forall K$, il suffit de le voir pour $K$ de la forme \deleted{\unsure{K}} $\mathbb E^i$. Or on a alors

% 21, (2)
\[
 F(\mathbb E^i)\longrightarrow\operatorname{Hom}^{\bullet}_{k}(\mathbb E_i^{\bullet},\deleted{\unsure{\mathbb E^{\bullet}}}F(C^{\bullet})) \tag*{(2)}
\]

% 22
où on a \unsure{évidemt|évident|évidemment}

% 23
\[
 \operatorname{Hom}^{\bullet}_{k}(\mathbb E_i^{\bullet},\underset{\substack{\rotatebox{-90}{$\in$}\\Ab^{\bullet}}}{\unsure{D^{\bullet}|\mathcal C^{\bullet}}})\simeq C^i
\]

% 24–25
et on voit de suite que moyennant cet iso, (2) est l'identité, cqfd

% archive, bottom left
% archive sheet number: 44
\scanpage{46}

% 1–3, heading and opening
\uline{Ex 3} Les \unsure{triples|triplets|théorie} de Dold-Puppe nous donnent une \deleted{\unsure{foncteur} équivalence de catégories} équivalence de catégories (pour tte catégorie \deleted{$\unsure{Ab_k|Ab_k^\bullet}\quad Ab_k$.} additive \unsure{avec} facteurs directs)

% 1–3, upper left margin
% slanted
\margin{\uline{NB} Le foncteur $DP$ est défini de \unsure{ple.|faç.} \unsure{fid.} pour des $B$ additives (\unsure{qui ont|karoubiennes} des facteurs directs)}

% 4, abbreviations in left margin
\margin{$ND=$ non dégénéré \quad $DP=$ Dold-Puppe \deleted{\unsure{de $Ab$}}}

% 4–5, equivalences
\diagram{\begin{tikzcd}[ampersand replacement=\&,column sep=small]
 B_{\bullet}\arrow[r,shift right=1.1ex,"{DP^{*}}"']\& B_{*}\arrow[l,bend right=35,"{ND_{\bullet}}"']\arrow[l,phantom,"\simeq"] \& \text{donc aussi}\ \substack{\text{(en appliquant}\\\text{à } B^\circ\text{)}}\&
 B^{\bullet}\arrow[r,shift right=1.1ex,"{DP_{*}}"']\& B^{*}\arrow[l,bend right=35,"{ND^{\bullet}}"']\arrow[l,phantom,"\simeq"]
\end{tikzcd}}

% 4–5, left margin under abbreviations
% slanted
\margin{équiv. compatibles aux foncteurs additifs}

% 6
en particulier

% 7
\[
 Ab_{k\,\bullet}\overset{\sim}{\longrightarrow}Ab_{k\,*},\qquad Ab_k^{\bullet}\overset{\sim}{\longrightarrow}Ab_k^{*}
\]

% 8–13
\deleted{Il est immédiat} \deleted{\unsure{Donc}.} l'image \unsure{essentielles|essentielle} de $\mathbb K_{\bullet}^{\deleted{\unsure{k}}}$ et $\mathbb K^{\bullet}_{\deleted{\unsure{k}}}$ par ces deux foncteurs sont notés $\mathbb K_{*}$ et $\mathbb K^{*}$. On a donc nos diagrammes d'équivalences canoniques

% 9–13, left margin
% slanted
\margin{$k$ est sous-entendu dans ces notations, sinon, c'est $\mathbb K_{!}(k)$ avec indices ou exposant $\bullet$ ou $*$.}

% 14–15, cancelled diagram
\diagram{\begin{tikzcd}[ampersand replacement=\&,execute at end picture={\draw (dlla.south west) -- (dula.north east) -- (dlla.south east) -- (dura.north west) -- (dlra.south west) -- (dura.east);}]
 |[alias=dula]|\deleted{Ab_{k\,\bullet}}\arrow[r,shift right,"{\deleted{DP}}"']\&|[alias=dura]|\deleted{Ab_{k\,*}}\arrow[l,shift right,"{\deleted{ND}}"']\\
 |[alias=dlla]|\deleted{Ab_k^{\bullet}}\arrow[r,shift right]\&|[alias=dlra]|\deleted{Ab_k^{*}}\arrow[l,shift right]
\end{tikzcd}}

% 16–20, five-object diagram
\diagram{\begin{tikzcd}[ampersand replacement=\&,row sep=2.4em,column sep=2.4em]
 \mathbb K_{\bullet}\arrow[rr,shift right,"{DP^{*}}"']\arrow[dd,leftrightarrow,"{\vee\,\circ}", "{\deleted{\unsure{\mathrm{can}}}}"'] \& \& \mathbb K_{*}\arrow[ll,shift right,"{ND_{\bullet}}"']\arrow[dd,leftrightarrow,"{\circ\,\vee}", "{\deleted{\unsure{\mathrm{can}}}}"']\\
 \&\mathbb K\arrow[ul,"{\alpha_{\bullet}}"]\arrow[ur,"{\alpha_{*}}"']\arrow[dl,"{\alpha^{\bullet}\ \circ}"]\arrow[dr,"{\alpha^{*}\ \circ}"']\&\\
 \mathbb K^{\bullet}\arrow[rr,shift right,"{DP_{*}}"']\&\&\mathbb K^{*}\arrow[ll,shift right,"{ND^{\bullet}}"']
\end{tikzcd}}

% 16–23, long marginal convention, left
\margin{Convention d'écriture. pour un objet $C$ de $K$, on note $C_{\bullet}, C_{*}, C^{\bullet}, C^{*}$ les objets correspondants obtenus en appliquant $\alpha_{\bullet},\alpha_{*},\alpha^{\bullet},\alpha^{*}$, et ses comp. sont notés $C_n,C_{[n]},C^n,C^{[n]}$ (Attention, \uline{ne pas} oublier le \uline{signe} $[\ ]$ et $\mathbb E_{\bullet}^{i},\mathbb E_{*}$ $\mathbb E^{\bullet},\mathbb E^{*}$, $\mathbb E^{i}$}

% 21–22
Les flèches \deleted{\unsure{sont}} marquées d'un signe $\circ$ sont des foncteurs \uline{contravariant}.

% 23–25
On a

\uline{Proposition 4} \deleted{\illegible{2}} On a $\forall i\geq0$ un isom canonique

% archive, bottom left
% archive sheet number: 45
\scanpage{47}

% 1
\[
 \alpha_{*}(\mathbb E^i)\simeq DP(C_{\bullet}^{i})\simeq\overset{i+1}{\bigwedge}\Phi_{*}
\]

% 2
\[
 (\text{d'où}\quad\alpha^{*}(\mathbb E^i)\simeq DP(C_{i}^{\bullet})\simeq\overset{i+1}{\bigwedge}\check{\Phi}_{*}\;)
\]

% 3–5, cancelled at upper edge
\deleted{L'\unsure{équ.|iso} déf. entre les $\alpha_{*}(\mathbb E^i)$ est l'\unsure{équiv. des} $T_{*}\wedge$ \unsure{que} le foncteur $ND:\mathbb K_{*}\longrightarrow\mathbb K_{\bullet}$ est \unsure{donné} \unsure{comme} de $(\overset{\unsure{\bullet\geq1|i+1}}{\bigwedge})\Phi_{*}$.}

% 5–7, crossed heading and formula
\deleted{Par suite \uline{Cor. 1} Le foncteur $DP:Ab_{k\,\bullet}\longrightarrow\unsure{Ab_{k\,*}}$}

% 8–11
Appliquant le \unsure{triple|théorie} de Dold-Puppe à $K$ \deleted{(qui est bien additive et stable par facteurs directs, du moins si $k$ est principal, p. ex $k=\mathbb Z$)} on trouve dans $K$ un complexe

% 12
\[
 \deleted{\unsure{C^{*}}}\qquad C^{*}=DP(C^{\bullet})
\]

% 13–14
\deleted{et \unsure{par|dans} Yoneda, pour} (qui est donc défini dans \deleted{\unsure{$K_k$}} $K(k)$ quel que soit $k$), d'où pour tte $B$ \inserted{$k$-}additive

% 15–18, triangle with intermediate restriction
\diagram{\begin{tikzcd}[ampersand replacement=\&,row sep=large]
 \underline{\operatorname{Hom}}_{k\,\mathrm{lin}}(\mathbb K,B)\arrow[d,"\mathrm{restr}","\wr"']\arrow[ddr,bend left=15,"{F\mapsto F(C^{*})}"]\&\\
 \underline{\operatorname{Hom}}_{k\text{-}\mathrm{lin}}(\unsure{\mathscr A|\mathcal A},B)\arrow[d,"{F\mapsto F(C^{\bullet})}","\wr"']\&\\
 B^{\bullet}\arrow[r,leftrightarrow,"{DP_{\unsure{\bullet|\,}}}"]\& B^{*}
\end{tikzcd}}

% 19–23
donc un foncteur $F\mapsto F(C^{*})$ qui est pl. fid. Si $B$ est karoubienne, $DP_B$ est une équivalence, donc aussi $F\mapsto F(C^{*})$. Donc $C^{*}$ apparaît comme objet \unsure{cosimplicial|simplicial} \uline{universel} dans cat. $k$-Additive karoubienne

% 24–25
\uline{Prop 5} Si $B$ est $k$-additive karoubienne, alors $F\mapsto F(C^{*})$ est une équivalence

% 26
\[
 \underline{\operatorname{Hom}}_{k\text{-}\mathrm{lin}}(\mathbb K,B)\longrightarrow B^{*}
\]

% 27–28
On a un isom canonique \inserted{bifonct.l} en $F\in\operatorname{Ob}\underline{\operatorname{Hom}}_{k\text{-}\mathrm{lin}}(\mathbb K,B)$ et $x\in\operatorname{Ob}\mathbb K$.

% 28, lower edge below Hom
\margin{\unsure{$\wr$}}

% archive, bottom left
% archive sheet number: 46
\scanpage{48}

% 1, isolated mark at upper right
\margin{\unsure{?|2}}

% archive, bottom left
% archive sheet number: 47
\scanpage{49}

% 1–4
Soit $\mathbb P_{\bullet}$ la catégorie des complexes de $k$-chaînes $L_{\bullet}$ qui sont libres de r. f. en chaque dim, et qui n'ont qu'un nb fini de composantes non nulles. Une suite

% 1–7, left margin
% slanted
\margin{Il vaut mieux construire $\mathbb P$ ``en l'air'', puis le ``réaliser'' \deleted{\unsure{quatre}} $\mathbb P_{\bullet},\mathbb P_{*},\mathbb P^{\bullet},\mathbb P^{*}$ comme pour $\mathbb K$.}

% 5
\[
 0\longrightarrow L'_{\bullet}\longrightarrow L_{\bullet}\longrightarrow L''_{\bullet}\longrightarrow 0
\]

% 6–8
\inserted{d'hom dans $\mathbb P_{\bullet}$} est dite exacte spéciale, si $\forall n\geq0$, la suite $0\to L'_n\longrightarrow L_n\longrightarrow L''_n\longrightarrow0$ est exacte. Soit $\Sigma_{\bullet}$ la classe de ces suites.

% 8–11
Ainsi $\mathbb P_{\bullet}$ est muni de $\mathbb E_{\bullet}^{\bullet}$ (complexe de cochaînes) et de $\Sigma_{\bullet}$. On va voir que pour ces données, il est sol. d'un pb universel.

% 11, left margin
\margin{!}

% 12–13
\deleted{Soit \unsure{$A$}} $B$ une cat. additive. Un ens. $\Sigma$ de suites

% 14, (x)
\[
 0\longrightarrow A'\xrightarrow{q} A\xrightarrow{p}A''\longrightarrow0 \tag*{(x)}
\]

% 15–16
dans $B$ est dit ``une classe de suites exactes'' si elle satisfait aux conditions

% 17–20
Sex 1) Si $(x)\in\Sigma$, alors $A'\overset{\sim}{\longrightarrow}\operatorname{Ker}p$, $\operatorname{Coker}q\overset{\sim}{\longrightarrow}A''$, i.e. les suites $0\longrightarrow A'\longrightarrow A\longrightarrow A''$ et $\deleted{\unsure{0\to}}A'\longrightarrow A\longrightarrow A''\longrightarrow0$ sont exactes (i.e. $\forall B\in\operatorname{Ob}B$, les suites de groupes ab.

% 21–22
\[
 \left\{\begin{aligned}
 0&\longrightarrow\operatorname{Hom}(B,A')\longrightarrow\operatorname{Hom}(B,A)\longrightarrow\operatorname{Hom}(B,A'')\\
 \deleted{\text{et}}\quad0&\longrightarrow\operatorname{Hom}(A'',B)\longrightarrow\operatorname{Hom}(A,B)\longrightarrow\operatorname{Hom}(A',B)
 \end{aligned}\right.
\]

% 23
sont exactes

% 24–25
Sex 2) $\Sigma$ stable par iso, par sommes finies et contient les suites de la forme

% 26–27
\[
 \begin{aligned}0&\longrightarrow A\xrightarrow{\mathrm{id}} A\longrightarrow0\longrightarrow0\\
 0&\longrightarrow0\longrightarrow A\xrightarrow{\mathrm{id}} A\longrightarrow0\end{aligned}
\]

% 28
(donc aussi les suites de la forme

% 29
\[
 0\longrightarrow A\xrightarrow{\mathrm{inj}}A\times B\xrightarrow{\mathrm{pr.}}B\longrightarrow0\qquad)
\]

% archive, bottom left
% archive sheet number: 48
\scanpage{50}

% 1–3
Un $\Sigma$-monomorphisme (épimorphisme) $A\longrightarrow B$ de $B$ est un morph. qui s'insère dans une suite $\in\Sigma$

% 4
\[
 0\longrightarrow A\longrightarrow B\longrightarrow K\longrightarrow0\quad(\text{resp.}\quad0\longrightarrow P\longrightarrow A\longrightarrow B\longrightarrow0)
\]

% 5–7
Un morphisme \inserted{$A\longrightarrow B$} \deleted{\unsure{est}} est \inserted{$\Sigma$-}strict, s'il a un noyau $P$ et un conoyau $Q$, \deleted{\unsure{et si|objet} $[\operatorname{Coim}\overset{\sim}{\longrightarrow}\operatorname{Im}$,} \inserted{et si $P\longrightarrow A$} $(B\longrightarrow Q)$ est un $\Sigma$-mono ($\Sigma$-épi),

% 5–7, outer left margin
% slanted
\margin{définir \unsure{$\Sigma$-mo.|$\Sigma$-simpl.} et \unsure{$\Sigma$-épi|$\Sigma$-co.} \deleted{\illegible{1}}}

% 7–9
Il revient au même de dire que l'on peut trouver des suites $\in\Sigma$

% 10–11
\[
 \begin{aligned}
 0&\longrightarrow P\longrightarrow A\longrightarrow R\longrightarrow0\\
 0&\longrightarrow R\longrightarrow B\longrightarrow Q\longrightarrow0
 \end{aligned}
\]

% 12–14
telles que $A\longrightarrow B$ soit le composé $A\to R\longrightarrow B$. (Il faut peut-être exiger qu'un composé de $\Sigma$-mono (resp $\Sigma$-épi) est itou.)

% 15–16
On définit de façon évidente les \inserted{objets} $\Sigma$-projectifs et $\Sigma$-injectifs. Une \deleted{\inserted{suite} complexe est dit} \deleted{complexe de cochaînes (disons)}

% 17–18, cancelled
\deleted{$\longrightarrow C^{i-1}\longrightarrow \overset{\deleted{\unsure{d}}}{C^i}\longrightarrow C^{i+1}\longrightarrow\cdots$ est dit spécial.}

% 19–22
\[
 A\xrightarrow{u}B\xrightarrow{v}C
\]
\deleted{est dite spéciale si \unsure{$vu=0$}, et} $\Sigma$-spéciale \deleted{\unsure{si}} si elle \deleted{elle \unsure{déf.} s'insère dans des suites (de $\Sigma$)} \deleted{$vu=0$, $u$ est spécial et s'insère ds une suite \unsure{$v=0$|$vu=0$}, \unsure{sa}} de suites exactes $\in\Sigma$

% 23–26, stepped sequences
\diagram{\begin{tikzcd}[ampersand replacement=\&,column sep=.45em,row sep=1.3em]
 0\arrow[r]\&\operatorname{Ker}u\arrow[r]\& A\arrow[r]\&\operatorname{Coim}u\arrow[r]\arrow[d,"\wr"]\&0\&\&\&\&\&\\
 \&\&0\arrow[r]\&\operatorname{Im}u\arrow[r]\&\operatorname{Ker}v\arrow[r]\arrow[d,equal,"\wr"]\&H\arrow[r]\&0\&\&\&\\
 \&\&\&0\arrow[r]\&\operatorname{Ker}v\arrow[r]\& B\arrow[r]\&\operatorname{Coim}v\arrow[r]\arrow[d,"\wr"]\&0\&\&\\
 \&\&\&\&\&\deleted{0}\arrow[r]\&\operatorname{Im}v\arrow[r]\& C\arrow[r]\&\unsure{\operatorname{Cok}v|\operatorname{Coker}v}\arrow[d]\&\\
 \&\&\&\&\&\&\&\&0\&
\end{tikzcd}}

% 27–32
(c'est une notion autoduale) On dit qu'un complexe $C^{\bullet}$ dans $B$ est \uline{$\Sigma$-spécial} si …, $\Sigma$-acyclique si de \uline{plus} les $H^i(C^{\bullet})$ sont nuls, une $\Sigma$-résolution si c'est un complexe de cochaînes ($C^i=0$ si $i<0$) et $H^i=0$ pour $i\deleted{\unsure{\geq}}\unsure{>|\neq|\geq}0$.

% archive, bottom left
% archive sheet number: 49
\scanpage{51}

% 1
Il faudrait prouver le

% 2, outer left
\margin{!}

% 2–6
\uline{Th} Soit $B$ une cat. \deleted{addit} $k$-additive avec famille $\Sigma$ de suites exactes, \deleted{soit} $\underline{\operatorname{Hom}}_{k,\Sigma\text{-}\mathrm{ex}}(\mathbb P_{\bullet},B)$ la \inserted{$k$-}catégorie \deleted{pleine de} $\underline{\operatorname{Hom}}_k(\mathbb P_{\bullet},B)$ formée des foncteurs ``$\Sigma$-exacts'' i.e. transformant suites de $\Sigma_{\bullet}$ en suites de $\Sigma$,

% 2–4, right marginal addition
\margin{\unsure{NB|où}, tte flèche a un $\Sigma$-noyau. \unsure{Gr. cas|p. ex.|i.e.} $B$ \uline{cat. abélienne} $\Sigma$ ttes les suites ex.}

% 7–10, overlaid drafts
\deleted{Considérons les foncteurs exacts} \deleted{Soit $\unsure{\underline{\operatorname{Hom}}_{\Sigma\text{-}\mathrm{ex}}(\mathbb E_{\bullet},B)|\underline{\operatorname{Hom}}_{\Sigma\text{-}\mathrm{ex}}(\mathbb P_{\bullet},B)}$ la sous-catégorie pleine} $\underline{\operatorname{Hom}}_{k,\Sigma\text{-}\mathrm{sex}}(\mathbb P_{\bullet},B)$ la sous-catégorie pleine de celles qui \deleted{transforment $\Sigma$-mono en $\Sigma$-mono}. \inserted{commutent \deleted{\unsure{à}} aux $\Sigma$-noyaux}

% 11
a) Le foncteur restriction $F\longmapsto F(\mathbb E^{\bullet}_{\bullet})$

% 12
\[
 \underline{\operatorname{Hom}}_{k,\Sigma\text{-}\mathrm{sex}}(\mathbb P_{\bullet},B)\longrightarrow\deleted{\underline{\operatorname{Hom}}_{k\,\mathrm{lin}}(\mathbb K,B)\simeq}B^{\bullet}
\]

% 13–17, revised sentence
\deleted{est une équivalence de catégories} \deleted{(pour les foncteurs} \deleted{est pleinement fidèle, et son $\underline{\operatorname{Hom}}_{k,\Sigma\text{-}\mathrm{sex}}$) image essentielle est formée des $C^{\bullet}\in\operatorname{Ob} B^{\bullet}$ tels que les $d^i:C^i\to C^{i+1}$ admettent des \inserted{$\Sigma$-}noyaux}

% 14–16, substitution in left margin
% slanted
\inserted{est une équivalence de cat.}

% 18–22
[donc \deleted{\unsure{$K$}} $(\mathbb P_{\bullet},\mathbb E^{\bullet},\Sigma_{\bullet})$ est universel pour les triples $(B,C^{\bullet},\Sigma)$ tels que les $d^i:C^i\longrightarrow C^{i+1}$ aient des $\Sigma$-noyaux, \unsure{qud|quand} on prend comme morphismes les foncteurs \unsure{convenmt} $\Sigma$-exacts à g.]

% 23
b) Le foncteur \deleted{restriction} \inserted{$F\longmapsto F(\mathbb E^{\bullet}_{\bullet})$.}

% 24
\[
 \underline{\operatorname{Hom}}_{k\text{-}\Sigma\text{-}\mathrm{ex}}(\mathbb P_{\bullet},B)\longrightarrow B^{\bullet}
\]

% 25–28
(qui est pl. fid. par a)) a comme image essentielle les $C^{\bullet}\in\operatorname{Ob}B^{\bullet}$ qui sont des $\Sigma$-résolutions (i.e. $\Sigma$-spéciaux et avec $H^i(C^{\bullet})=0$ si $i>0$).

% 28–31, bracketed cancellation at foot
\deleted{[donc le triple $(\mathbb P_{\bullet},\mathbb E_{\bullet}^{\bullet},\Sigma_{\bullet})$ apparaît comme universel parmi les $(B,C^{\bullet},\Sigma)$ tels que $C^{\bullet}$ soit \deleted{\unsure{exacte|construite}} une $\Sigma$-résolution, -- les}

% archive, bottom left
% archive sheet number: 50
\scanpage{52}

\deleted{
% 1–2
foncteurs admissibles étant maintenant les foncteurs $\Sigma$-exacts…)

% 3
\uline{Indications sur dém}

% 4–7
Dans $\mathbb P_{\bullet}$, il est évident que \uline{$\mathbb E_{\bullet}^{\bullet}$ est une $\Sigma_{\bullet}$-spéciale résolution}. De façon précise, si $\deleted{\unsure{-|\to}}\Psi_{\bullet}^{i}=\deleted{\unsure{Z^{i}(\mathbb E_{\bullet})|H^i(C^\bullet)}} k[\deleted{\unsure{\,i\,}}i]$ placé en degré $i$ on voit que l'on a des suites $\in\Sigma_{\bullet}$

% 8–10, staggered rows
\[
 \begin{aligned}
 &0\longrightarrow\Psi_{\bullet}^{0}\longrightarrow C_{\bullet}^{0}\longrightarrow\Psi_{\bullet}^{1}\longrightarrow0\\
 &\hspace{4em}0\longrightarrow\Psi_{\bullet}^{1}\longrightarrow C_{\bullet}^{1}\longrightarrow\Psi_{\bullet}^{2}\longrightarrow0\\
 &\hspace{8em}0\longrightarrow\Psi_{\bullet}^{2}\longrightarrow C_{\bullet}^{2}\longrightarrow\Psi_{\bullet}^{3}\longrightarrow\unsure{0|\cdots}
 \end{aligned}
\]

% 11–14
De plus, on voit que pour $L_{\bullet}\in\operatorname{Ob}\unsure{\overline{\mathbb P}_{\bullet}|\mathbb P_{\bullet}|\mathcal P_{\bullet}}$, $L_{\bullet}$ peut s'écrire comme extension successive d'objets de la forme $(\Psi_{\bullet}^{i})^{\unsure{\rho_i|f_i}}$ pour la filtration croissante

% 15–16
\[
 \mathrm{Fil}_i(L_{\bullet})_j=\begin{cases}L_j&\text{si }j\leq i\\0&\text{si }j>i\end{cases}
\]

% 17
on a un isom can

% 18–19
\[
 \begin{aligned}
 \deleted{\mathrm{Fil}}\ \mathrm{Gr}_i(L_{\bullet})&\simeq L_i\ \text{placé en degré }i\\
 &\underset{\mathrm{can}}{\simeq}\Psi_{\bullet}^{i}\otimes L_i\underset{\text{non can}}{\simeq}(\Psi_{\bullet}^{i})^{\rho_i}
 \end{aligned}
\]

% 20
\uline{Enfin}, les $\deleted{\unsure{\mathbb E}}C_{\bullet}^{i}$ sont $\Sigma_{\bullet}$-projectifs

% 21–22
Soit $B^{\bullet}_{\deleted{\unsure{\Sigma}}\Sigma\text{-}\mathrm{sex}}$ (resp. $B^{\bullet}_{\Sigma\text{-}\mathrm{ex}}$) la sous-catégorie pleine de $B^{\bullet}$ décrite dans a) (resp. b)).

% 23
On va définir des foncteurs

% 24–25, square
\diagram{\begin{tikzcd}[ampersand replacement=\&]
 B^{\bullet}_{\Sigma\text{-}\mathrm{sex}}\arrow[r]\arrow[d]\&\underline{\operatorname{Hom}}_{\Sigma\text{-}\mathrm{sex}}(\mathbb P_{\bullet},B)\arrow[d]\\
 B^{\bullet}_{\Sigma\text{-}\mathrm{ex}}\arrow[r]\&\underline{\operatorname{Hom}}_{\Sigma\text{-}\mathrm{ex}}(\mathbb P_{\bullet},B)
\end{tikzcd}}

% 26
(le premier induisant le deuxième,
}

% archive, bottom left
% archive sheet number: 51
\scanpage{53}

% 1, cancelled
\deleted{Soit $\mathbb P\supset\mathbb K$ une \unsure{sur-catégorie|sous-catégorie} $k$-additive.}

% 2, cancelled upper pair
\diagram{\begin{tikzcd}[ampersand replacement=\&,execute at end picture={\draw (current bounding box.south west) -- (current bounding box.north east);}]
 \deleted{\underline{\operatorname{Hom}}_{k\text{-}\mathrm{lin}}(\mathbb P,B)}\arrow[r,shift left,"{\deleted{\rho}}"]\&\deleted{B^{\bullet}}\arrow[l,shift left,"{\deleted{\sigma}}"]
\end{tikzcd}}

% 2–4, upper right
\margin{\deleted{$B$ cat. additive avec $\Sigma$ \deleted{\unsure{de}} classe de suites exactes et $\Sigma$-noyaux}}

% 3, left, alongside triangle
\deleted{\unsure{pl. f.}}

% 3–4, cancelled triangle
\diagram{\begin{tikzcd}[ampersand replacement=\&,execute at end picture={\draw (current bounding box.south west) -- (current bounding box.north east);}]
 \deleted{\underline{\operatorname{Hom}}_{k\text{-}\mathrm{lin}}(\mathbb P,\widehat B)}\arrow[d]\&\\
 \deleted{\underline{\operatorname{Hom}}_{k\text{-}\mathrm{lin}}(\mathbb P,B)}\arrow[r,"{\deleted{\rho}}"]\&\deleted{B^{\bullet}}\arrow[ul,"{\deleted{\sigma}}"]
\end{tikzcd}}

% 5–6
$B$ catégorie additive, \deleted{avec noyaux de flèches,} \deleted{\inserted{classe $\Sigma$ de suites exactes}} où tt morphisme admet un $\Sigma$-noyau

% 7, new pair of functors
\diagram{\begin{tikzcd}[ampersand replacement=\&]
 \underline{\operatorname{Hom}}_{k,\unsure{\Sigma\text{-}\mathrm{sex}}}(\mathbb P_{\bullet},B)\arrow[r,shift left,"\rho"]\& B^{\bullet}\arrow[l,shift left,"{\unsure{\sigma|\substack{\sigma\\\mathrm{adj}}}}"]
\end{tikzcd}}

% 7, large note at right
% slanted
\margin{non, il faut condition de platitude ?}

% 8–11, below left Hom
cat. des foncteurs $k$-lin $F:\mathbb P\longrightarrow B$ qui transforment suites $\in\Sigma$ en suites exactes à \unsure{d|g}

% 12–13, defining brace
\[
 \left\{\begin{aligned}
 \rho(F)&=F(\mathbb E^{\bullet})\\
 \sigma(C^{\bullet})(L_{\bullet})&\simeq\operatorname{Hom}^{\bullet}_{k}(L_{\bullet}^{\vee},C^{\bullet})\left[=\operatorname{Hom}^{\bullet}_{k}(\operatorname{Hom}(L_{\bullet},\mathbb E^{\bullet}_{\deleted{\unsure{\bullet}}}),C^{\bullet})\right]
 \end{aligned}\right.
\]

% 13, below the second member
\margin{défini grâce à l'existence des $\Sigma$-noyaux}

% 14
Formule d'adjonction

% 15–16, with vertical identification
\diagram{\begin{tikzcd}[ampersand replacement=\&]
 \operatorname{Hom}(\rho(F),C^{\bullet})\arrow[r,"\sim"]\arrow[d,phantom,sloped,"\simeq"]\&\operatorname{Hom}(F,\sigma(C^{\bullet}))\\
 \operatorname{Hom}(F(\mathbb E^{\bullet}),C^{\bullet})\&
\end{tikzcd}}

% 17
Le deuxième mb est formé des hom fonctoriels en $L_{\bullet}$

% 18
\[
 F(L_{\bullet})\longrightarrow\operatorname{Hom}^{\bullet}_{k}(\underset{\substack{\Vert\\L^{\bullet}}}{L_{\bullet}^{\vee}},C^{\bullet})
\]

% 19
i.e. des hom ``fonctoriels en $L_{\bullet}$''

% 20
\[
 L^{\deleted{\unsure{k}}\bullet}\otimes_k F(L_{\bullet})\longrightarrow C^{\bullet}
\]

% 21
Or on a $L^i\simeq\operatorname{Hom}(L_{\bullet},\mathbb E^i)$

% 22
\[
 F(L)\otimes L^i\deleted{\otimes F(L_{\bullet})}\simeq\operatorname{Hom}(L_{\bullet},\mathbb E^i)\otimes F(L_{\bullet})\xrightarrow{\mathrm{taut.}}F(\mathbb E^i)
\]

% 23–24
et on trouve qu'on a ainsi la \deleted{``fonctorialité} compatibilité à l'op. diff. en $L^{\bullet}$ et $C^{\bullet}$, et ``fonctorialité

% archive, bottom left
% archive sheet number: 52
\scanpage{54}

% 1
en $L_{\bullet}$'', ce qui donne si on veut

% 2
\[
 F\xrightarrow{u}\sigma\rho F
\]

% 3–4
hom fonctoriel en $F$, \inserted{i.e. un} \deleted{\unsure{d'où}} \deleted{\unsure{D}} hom d'adjonction partiel

% 5–6, triangular comparison
\diagram{\begin{tikzcd}[ampersand replacement=\&,column sep=1.6em]
 \operatorname{Hom}(\rho F,C^{\bullet})\arrow[rr]\arrow[dr,dashed]\&\&\operatorname{Hom}(F,\sigma C^{\bullet})\\
 \&\operatorname{Hom}(\sigma\rho F,\sigma C^{\bullet})\arrow[ur,dashed,"{\alpha\mapsto\alpha\circ u}"']\&
\end{tikzcd}}

% 7–8
à prouver que c'est un iso -- vérification facile.

% 9
L'hom can

% 10
\[
 \rho\sigma C^{\bullet}\longrightarrow C^{\bullet}
\]

% 11
est un iso, i.e. $\sigma$ est pl. fid. En effet,

% 12–14, with identification below
\[
 \rho(\sigma(C^{\bullet}))=\sigma(C^{\bullet})(\mathbb E^{\bullet})=
 \bigl\{\underbrace{\sigma(C^{\bullet})(\mathbb E_{\bullet}^{i})}_{\substack{\operatorname{Hom}_{k}^{\bullet}(\mathbb E_i^{\bullet},C^{\bullet})\\ \Vert\\ C^{i}}}\bigr\}_{i}
 \simeq(C^{i})_{i}\quad\text{OK.}
\]

% 15
Ainsi, on a une inclusion \unsure{pl.|plt} fidèle

% 16
\[
 B^{\bullet}\overset{\sigma}{\lhook\joinrel\longrightarrow}\underline{\operatorname{Hom}}_{k,\Sigma\text{-}\mathrm{sex}}\ \subset\ \underline{\operatorname{Hom}}_{k\text{-}\mathrm{lin}}
\]

% 17–18
où $\Sigma$-sex désigne les foncteurs qui transf. suite de $\Sigma$ en suite exacte à g.

% 19–23
Quelle est l'image essentielle de $B^{\bullet}$ ? (On donnera des exemples tantôt de foncteurs ds l'image ess.) est-ce tout $\underline{\operatorname{Hom}}_{k,\Sigma\text{-}\mathrm{sex}}$ ? Elle est formée des $F$ telles que

% 24
\[
 F\xrightarrow{u}\sigma\rho(F)
\]

% 25
soit un iso. Or on a

% archive, bottom left
% archive sheet number: 53
\scanpage{55}

% 1–2
\[
 F(\mathbb E^i_{\bullet})\longrightarrow\sigma(\rho(F))(\mathbb E^i_{\bullet})=\operatorname{Hom}^{\bullet}_{k}(\mathbb E^{\deleted{\unsure{i}}\bullet}_{i},\underset{\substack{\Vert\\F(\mathbb E_{\bullet}^{\bullet})}}{\rho(F)})\simeq F(\mathbb E^i_{\bullet})
\]

% 3–6
\[
 \begin{aligned}
 F(\Psi^{i}_{\bullet})&\longrightarrow\operatorname{Hom}^{\bullet}_{k}(\Psi^{\deleted{\unsure{i}}\bullet}_{i},F(\mathbb E^{\bullet}))\simeq Z^{i}(F(\mathbb E^{\bullet}))\\
 &\simeq\operatorname{Ker}(F(\mathbb E^i)\to F(\mathbb E^{i+1}))\\
 &\simeq F(\underbrace{\operatorname{Ker}(\mathbb E^i\longrightarrow\mathbb E^{i+1})}_{\Sigma\text{-}\operatorname{Ker}\,!})=F(\Psi^i)
 \end{aligned}
\]

% 7–13
Donc on voit que $u(L_{\bullet}):F(L_{\bullet})\to\sigma\rho(F)(L_{\bullet})$ est un iso sur les $\mathbb E^i$ et les $\Psi^{i}_{\bullet}$, donc sur \deleted{la} la cat. additive $\mathbb P_0\supset\mathbb K$ des sommes de tels.) Pour avoir $\forall L_{\bullet}$, on \inserted{essaye en} raisonnant par récurrence sur la longueur $n$ \deleted{\unsure{de $L_{\bullet}$|d'un $L_{\bullet}$}}, \inserted{de l'intervalle $[p,q]$ tel que $L^i=0$ si $i\notin[p,q]$.} C'est vu si $n=0$, ($L_{\bullet}$ somme \deleted{\unsure{de $\Psi_{\bullet}$}} \inserted{$\simeq(\Psi^{i}_{\bullet})^{r}$}), on suppose $n>0$ et supposons que c'est fait pour $n-1$.

% 14–17, truncation diagram; brace under the truncated terms
\diagram{\begin{tikzcd}[ampersand replacement=\&,column sep=.55em,execute at end picture={\path (trleft.south west) -- (trright.south east) coordinate[midway] (trmiddle); \draw (trleft.south west) .. controls +(0,-.4em) and +(-.8em,.3em) .. ([yshift=-.65em]trmiddle) node[below] {$\text{tronqué }L'_{\bullet}$} .. controls +(.8em,.3em) and +(0,-.4em) .. (trright.south east);}]
 \&\&\&\cdots\arrow[r]\&0\arrow[r]\&L_p\arrow[r,"\mathrm{id}"]\&L_p\arrow[r]\&0\&\& \left.\vphantom{\begin{matrix}x\\x\end{matrix}}\right\}\ (\mathbb E_{\bullet}^{p\deleted{\unsure{+1}}})^{r}\deleted{\unsure{L_{\bullet}''}}\\
 \cdots\arrow[r]\&0\arrow[r]\&|[alias=trleft]|L_{p+n}\arrow[r]\&L_{p+n-1}\arrow[r]\&\cdots\arrow[r]\&|[alias=trright]|L_{p+1}\arrow[u,"{d_{p+1}}"]\arrow[r]\&L_p\arrow[u,"{\mathrm{id}_{L_p}}"]\arrow[r]\&\deleted{\unsure{0}}\arrow[r]\&0\& L_{\bullet}\arrow[u]
\end{tikzcd}}

% 18
on trouve un \unsure{mor|m} de suites exactes (où $L_p\simeq k^r$)

% 19–20, square with exact rows
\diagram{\begin{tikzcd}[ampersand replacement=\&,column sep=1.2em]
 0\arrow[r]\&(\Psi^p_{\bullet})^r\arrow[r]\arrow[d,equal]\&L_{\bullet}\arrow[r]\arrow[d]\&\deleted{\unsure{\Psi}}L_{\bullet}^{\unsure{\prime|\prime\prime}}\arrow[r]\arrow[d]\&0\\
 0\arrow[r]\&(\Psi^p_{\bullet})^r\arrow[r]\&(\mathbb E_{\bullet}^{p})^r\arrow[r]\&(\Psi_{\bullet}^{p+1})^r\arrow[r]\&0
\end{tikzcd}}

% 19–24, starts to right of square
qui exprime $L_{\bullet}$ comme produit fibré de $L'_{\bullet}$ et $(\mathbb E_{\bullet}^{p})^r$ sur $(\Psi_{\bullet}^{p+1})^r$ (i.e. comme noyau d'un hom \deleted{surj.} $\Sigma$-\uline{épi} $(\mathbb E_{\bullet}^{p})^{r}\times L_{\bullet}^{\unsure{1|\prime}}\longrightarrow(\Psi_{\bullet}^{p+1})^{r}$). Donc on conclut par

% archive, bottom left
% archive sheet number: 54
\scanpage{56}

% 1
l'hyp de récurrence. On trouve donc

% 2
\[
 \unsure{\prime}\underline{\operatorname{Hom}}_{k,\Sigma\text{-}\mathrm{sex}}(\mathbb P_{\bullet},B)\overset{\approx}{\longrightarrow} B^{\bullet}
\]

% 3–5
(si $B$ muni d'une \deleted{catégorie} \inserted{classe $\Sigma$} de suites exactes telle que tt morphisme de $B$ admette un $\Sigma$-noyau)

% 6–10
\uline{NB} Si on \inserted{ne} suppose \deleted{\unsure{seulmt|seulement} que $B$ admette une classe $\Sigma$ de suites exactes, mais} pas même que tt morphisme de $B$ admette un noyau, l'argument donné montre \unsure{seulmt|seulement} que le foncteur évident $F\longmapsto F(\mathbb E_{\bullet}^{\bullet})$

% 11
\[
 \underline{\operatorname{Hom}}_{k,\Sigma\text{-}\mathrm{sex}}(\mathbb P_{\bullet},B)\underset{\rho}{\overset{\approx}{\lhook\joinrel\longrightarrow}}B^{\bullet}
\]

% 12–14
est pl. fid (avec adjoint $\sigma$ comme dessus, mais à valeurs dans $\underline{\operatorname{Hom}}_{k,\mathrm{lin}\deleted{\unsure{\mathrm{ex}|\Sigma\text{-}\mathrm{sex}}}}(\mathbb P,\widehat B)$ ; on a néanmoins $\sigma\rho\simeq\mathrm{id}$ i.e.

% 15–16
\[
 \deleted{F(L_{\bullet})\simeq}\quad F(L_{\bullet})\simeq\operatorname{Hom}^{\bullet}_{k}(L^{\bullet},\underbrace{F(\mathbb E^{\bullet})}_{\substack{\Vert\\C^{\bullet}}})
\]

% 17
pour tt $F\in\underline{\operatorname{Hom}}_{k,\Sigma\text{-}\mathrm{sex}}$.

% 18–23
\uline{Re NB} On \deleted{\unsure{est}} vérifie facilt, sous les conditions du NB, que si $F\longmapsto C^{\bullet}=F(\mathbb E^{\bullet})$ se correspondent $(F\in\operatorname{Ob}\underline{\operatorname{Hom}}_{k,\Sigma\text{-}\mathrm{sex}})$, alors $F$ est $\Sigma_{\bullet}$-exact ssi $C^{\bullet}$ est une $\Sigma$-résolution. (du moins si $B$ cat. abélienne avec $\Sigma$ ttes les suites exactes…)

% 24–26
\uline{Ex} Soit $\lambda_{\bullet}\in Ab_{k\,\bullet}$, considérons le foncteur $L_{\bullet}\longmapsto\operatorname{Hom}_{Ab_{k\,\bullet}}(\lambda_{\bullet},L_{\bullet})$, \inserted{$\mathbb P\longrightarrow Ab_k$} Il est $\Sigma$-exact à g (et \={m} exact \uline{à g}) -- en fait

% 26, lower right
\margin{TSVP}

% archive, bottom left
% archive sheet number: 55
\scanpage{57}

% 1
la vérification de la formule

% 2
\[
 \operatorname{Hom}(\lambda_{\bullet},L_{\bullet})\overset{\sim}{\longrightarrow}\operatorname{Hom}(L_{\bullet}^{\vee},\lambda_{\bullet}^{\vee})
\]

% 3
est évidente --

% archive, bottom left
% archive sheet number: 56
\scanpage{58}

% 1–2, heading
Construction intrinsèque de $P$ ($k$-lin)

$P\supset K$, et satisfait les conditions

% 3–5, upper part, a)
a) Le foncteur covariant $L\longmapsto\operatorname{Hom}(L,\mathbb E^0)^\vee$ est représentable par un \inserted{objet} $\Psi^0$, Considérons l'hom $\Psi^0\longrightarrow\mathbb E^0$ transposé de

% 6–7, upper part
\[
\begin{aligned}
\operatorname{Hom}(\Psi^0,L)&\longleftarrow\operatorname{Hom}(\mathbb E^0,L)\quad\bigl[\longleftarrow\operatorname{Hom}(\mathbb E^1,L)\\
\rotatebox{90}{$\simeq$}\quad&\hspace{5em}?\\
\operatorname{Hom}(L,\mathbb E^0)^\vee&
\end{aligned}
\]

% 8–10, pairing
déduit de l'accouplement
\[
\begin{aligned}
\operatorname{Hom}(\mathbb E^0,L)\times\operatorname{Hom}(L,\mathbb E^0)&\longrightarrow\operatorname{Hom}(\mathbb E^0,\mathbb E^0)\simeq k\\
(u,\ v)&\longmapsto vu
\end{aligned}
\]

% 11–14, left of argument
\diagram{\begin{tikzcd}[ampersand replacement=\&,column sep=small,row sep=small]
\&\mathbb E^0\arrow[dl,"d"']\arrow[dr,"u"]\&\&\\
\mathbb E^1\arrow[rr,"u'"']\&\&L\arrow[dr,"v"]\&\\
\&\&\&\mathbb E^0
\end{tikzcd}}

% 11–17, centre
Or soit $u'\colon\mathbb E^1\to L$, considérons $u=u'\circ d^0$, et $v\circ u=v\circ(u'\circ d)$, \inserted{$=(v\circ u')\circ d$} comme $vu'=0$, on a donc $v\circ u=0$, donc le composé \deleted{$\operatorname{Hom}$}
\[
\operatorname{Hom}(\Psi^0,L)\longleftarrow\operatorname{Hom}(\mathbb E^0,L)\longleftarrow\operatorname{Hom}(\mathbb E^1,L)
\]
est nul ce qui signifie aussi que
\[
\Psi^0\longrightarrow\mathbb E^0\longrightarrow\mathbb E^1
\]
est nul, donc que $\mathbb E^{\unsure{*|\bullet}}$ est $\Psi^0$-augmenté.

% 16–19, left margin
% slanted
\margin{foncteur $P\to$ \unsure{par|posé|pour} $L\longmapsto\operatorname{Hom}(L,\unsure{\varphi_{\bullet}^{*}|\varphi_{\bullet}^{0}})$}

% 18–19, across middle
Soit $\unsure{\varphi^\bullet|\Phi^\bullet}$ le complexe de cochaînes qu'il forme.

% 19–22, cancelled lines
\deleted{Notons aussi que} \deleted{Soit $P_0$ la sous-catégorie}

\deleted{b) La famille des $\mathbb E^i$ et $\Psi^0$ est génératrice stricte}

\deleted{c) $\varphi^\bullet$ est une résolution de}

% 23–26, above b)
$k$-lin. \deleted{engendrée} formée des sommes finies d'objets de la forme $\mathbb E^i,\varphi_{\bullet}^{0}$, i.e. les expressions de la forme
\[
L+(\Psi^0)^r\qquad r\in\mathbb N,\ L^{\unsure{\,|'}}\in\operatorname{Ob} K
\]

% 27, at left and in margin
\inserted{On a}
% slanted
\margin{plus bas \unsure{voir|si} \unsure{ce qu'il} \unsure{faudrait}}

% 27–30, rewritten b)
b) Les suites $0\to L'\to L$ \deleted{\unsure{telles que}} \inserted{$0\to L'\to L\to L''$,} \unsure{$0\to L'\to L\to L''\to0,\ L'\to L\to L''\to0$|$\to L''\to0,\ L'\to L\to L''\to0$}, telles que les $P\in P_0$ sont projectifs \unsure{pour celles|pour celles-ci|fournit}, i.e. telles que les $0\to L'\to L\to L''\to0$ soient exactes pour \unsure{les dits|les dits $P$}, sont exactes.

% 31–33, foot
\deleted{\unsure{c)}} Soit $\Sigma$ la classe de ces suites.

\deleted{\unsure{b)}}c) Le complexe augmenté $(\mathbb E^\bullet,\Psi^0,\varepsilon)$ est une $\Sigma$-résolution, i.e. $\varphi^\bullet$ est $\Sigma$-acyclique

% archive, bottom left
% archive sheet number: 57
